% META: {"id":"1SPE-VARALEA-CO-003","chapitre":"1SPE-VARIABLES-ALEATOIRES","type_objet":"corrige","exercice_ref":"1SPE-VARALEA-EX-003","capacites_codes":["C1"],"sources_inspiration":[],"status":"approved","origin":"generated","status_history":["generated","audited","accepted","approved"],"release_acceptance":"RELEASE_OWNER_BATCH_ACCEPTANCE","acceptance_closure_digest":"sha256:9b3ccf9a81c5520fbb7e03b7b2d3e7bf2057a02834b6d908bf3be5f49f3b553f","acceptance_receipt_digest":"sha256:4fad86f0282e0fa4e0419cca1ad6379ae7af5a280c04a4a90880f90a1c044242"}
% BEGIN-VERIFY
% from sympy import Rational
% p_r = Rational(3,6); p_b = Rational(2,6); p_v = Rational(1,6)
% assert p_r + p_b + p_v == 1
% assert p_b + p_v == Rational(1,2)
% END-VERIFY
\begin{corrige}{1SPE-VARALEA-EX-003}

\textbf{1.} $X$ prend les valeurs $1$, $3$ et $5$.

\medskip
\textbf{2.} Il y a $6$ boules au total.
\[
\begin{array}{|c|c|c|c|}
\hline
x_i & 1 & 3 & 5 \\
\hline
P(X=x_i) & \frac{3}{6} = \frac{1}{2} & \frac{2}{6} = \frac{1}{3} & \frac{1}{6} \\
\hline
\end{array}
\]
Vérification : $\frac{1}{2} + \frac{1}{3} + \frac{1}{6} = \frac{3+2+1}{6} = 1$. \checkmark

\medskip
\textbf{3.} $P(X \geq 3) = \frac{1}{3} + \frac{1}{6} = \frac{2+1}{6} = \frac{1}{2}$.

\end{corrige}
